% Passes 5798--5799: heavy-line disjointness Reye copy and W9 partial isometry.
\subsubsection{Transpose carries point--line Reye incidence to heavy--line disjointness}
\label{sec:pass5798-5799-reye-disjointness-duality}

The affine matrix coordinates sharpen the point/heavy outer twist to a full
incidence theorem.  Write the sixteen Reye lines as graphs
\[
 L_M=\{(w,Mw):0\ne w\in W\},
 \qquad M\in\operatorname{Hom}(W,V),
\]
and the twelve heavy supports as
\[
 H_{\phi,\psi}=\{(w,x):\phi(x)+\psi(w)=1\}.
\]
Restricting the heavy equation to $L_M$ gives the linear functional
$(\phi M+\psi)(w)$.  A nonzero functional on $\mathbb F_2^2$ takes value one
on exactly two of the three nonzero vectors, while the zero functional never
does.  Therefore
\[
 \boxed{
 |L_M\cap H_{\phi,\psi}|=
 \begin{cases}
 0,&\psi=\phi M,\\
 2,&\psi\ne\phi M.
 \end{cases}}
\]
The complete line--heavy intersection census is
\[
 \boxed{0^{48}\oplus2^{144}}.
\]

Define incidence by disjointness,
\[
 D_{M,(\phi,\psi)}=1
 \iff L_M\cap H_{\phi,\psi}=\varnothing.
\]
Every line is incident with three heavy blocks and every heavy block with four
lines.  Hence $D$ is a $12_4,16_3$ incidence structure.  In fact it is exactly
a second Reye copy.  Under
\[
 F(w,x)=(w^T,x^T),
 \qquad M\mapsto M^T,
\]
the original point--line relation $x=Mw$ becomes
\[
 x^T=w^TM^T,
\]
which is precisely the heavy-disjointness condition.  Thus
\[
 \boxed{D[M,h]=R[F^{-1}(h),M^T].}
\]
Transpose therefore carries the original point--line Reye configuration to the
heavy--line disjointness Reye configuration on the same sixteen-line family.
The carrier outer involution is an incidence dualization, not only a character
or stabilizer-class equivalence.

There is also a canonical linear bridge between the two common nine-spaces.
Let
\[
 C_R=4R-J_{12\times16},
 \qquad
 C_H=2H-J_{12}
\]
be the integer-centered point--line and point--heavy incidence matrices.  Then
\[
 \boxed{
 B=\frac14 C_R^TC_H=J_{16\times12}-4D.
 }
\]
Every entry of $B$ is $1$ or $-3$, its row and column sums vanish, and
\[
 \boxed{\operatorname{rank}_{\mathbb Q}B=9}.
\]
Moreover
\[
 \boxed{BB^T=C_R^TC_R},
 \qquad
 \boxed{B^TB=4C_H^TC_H},
\]
with
\[
 (C_R^TC_R)^2=64(C_R^TC_R),
 \qquad
 (C_H^TC_H)^2=16(C_H^TC_H).
\]
Consequently
\[
 \boxed{U=\frac18B}
\]
is a rank-nine partial isometry satisfying
\[
 \boxed{
 UU^T=E_{W_9,L},
 \qquad
 U^TU=E_{W_9,H}.
 }
\]
Thus the disjointness Reye matrix is itself the signed cross-transform that
identifies the heavy and line realizations of the common irreducible $W_9$.

\paragraph{Boundary.}
This is exact finite incidence geometry and rational linear algebra.  It does
not supply a quantum, continuum, dynamical, gauge, or particle-physics
identification.
